\section{Motivation}

The transformation of modern electricity systems is inseparable from the rise of variable renewable energy sources, especially wind and solar generation. As renewable penetration increases, the physical and economic behaviour of power systems becomes more strongly coupled to exogenous drivers such as solar radiation, wind speed, temperature, seasonality, and calendar effects. In consequence, electricity markets and grid operations are no longer well described by relatively stable, weakly perturbed processes. They instead exhibit pronounced non-stationarity, regime changes, heteroscedasticity, and cross-variable dependencies across price, load, renewable generation, and weather covariates \cite{hong2016probabilistic,weron2014electricity}. From a monitoring perspective, this means that ``normal'' system behaviour is now context-dependent rather than fixed.

This shift creates a direct challenge for anomaly detection. In traditional industrial settings, abnormality is often approximated through static thresholds, linear residual models, or control-chart style assumptions about stable operating distributions. However, such assumptions become fragile in weather-dependent energy systems. A price spike, a renewable shortfall, or a load deviation may be significant in one operating regime but entirely expected in another. As multiple surveys have emphasized, time-series anomaly detection is fundamentally harder than i.i.d. anomaly detection because temporal dependence, seasonality, contextual effects, and concept drift make the definition of abnormality intrinsically dynamic \cite{chandola2009anomaly,blazquez2021review}. The problem is amplified in electricity applications, where exogenous weather inputs and market responses jointly shape the observed trajectories \cite{hong2016probabilistic,weron2014electricity}.


Many practically important anomalies are also not simple point outliers. They may
be contextual, where a value is abnormal only relative to a particular weather or
calendar regime, or collective, where an extended temporal pattern is abnormal
even though no individual observation is extreme
\cite{chandola2009anomaly,blazquez2021review}. For example, depressed photovoltaic
output at noon on a clear summer day is suspicious, whereas the same output level
at night is expected. Likewise, a price increase may become operationally important
only when considered together with load, renewable availability, or other system
conditions. The practical task is therefore to identify meaningful departures from
expected multivariate behaviour rather than merely numerically extreme observations.

These observations motivate the central premise of this thesis: anomaly detection
for electricity markets and grid operations should be treated as a context-aware
forecasting and reasoning problem rather than as static outlier screening. This
formulation is well suited to weakly labelled operational settings, where trustworthy
anomaly labels are sparse or unavailable but historical records of system behaviour
are abundant \cite{chandola2009anomaly,shafer2008tutorial}. The specific learning,
thresholding, and reasoning components used to realise this premise are introduced
below and developed in the subsequent chapters.



\section{The Neuro-Symbolic Paradigm}

The methodological foundation of this work is the neuro-symbolic paradigm, which combines the pattern-learning ability of neural models with the transparency and constraint-enforcement capability of symbolic reasoning \cite{garcez2009neural,gebser2012answer,gebser2019multi}. Neural methods are effective at extracting nonlinear structure from noisy, high-dimensional observations, but they do not automatically respect explicit domain rules. Symbolic methods provide a natural language for expressing physical constraints, temporal consistency requirements, and operational logic, but they are not designed to discover such structure directly from raw continuous data. The two paradigms therefore address complementary parts of the anomaly-detection problem.

In the proposed framework, the neural component estimates expected system behaviour.
Given an engineered multivariate feature vector
$\mathbf{x}_t \in \mathbb{R}^{d}$ at time $t$, the forecasting backbone produces
an estimate of the target vector $\hat{\mathbf{y}}_t \in \mathbb{R}^{m}$,
\begin{equation}
\hat{\mathbf{y}}_t = f_{\theta}(\mathbf{x}_t),
\label{eq:intro_forecast}
\end{equation}
where $f_{\theta}(\cdot)$ is implemented as a lightweight Tiny discrete Cellular
Neural Network (dCeNN) encoder followed by an Extreme Learning Machine
(ELM)-style ridge regression head. These components are selected to support
structured representation learning, efficient training, and compact inference
\cite{chua1988cellular,huang2006extreme,hoerl1970ridge}.

After forecasting, the residual vector
\begin{equation}
\mathbf{r}_t = \mathbf{y}_t - \hat{\mathbf{y}}_t
\label{eq:intro_residual}
\end{equation}
quantifies deviation from expected behaviour. Residual thresholding first produces
statistical anomaly candidates, after which symbolic rules evaluate their contextual
and domain-level plausibility. This second stage allows information such as solar
conditions, calendar effects, persistence, and relationships among signals to
influence the final decision
\cite{gelfond1988stable,gebser2012answer,gebser2019multi}.

The resulting neuro-symbolic decision can be written schematically as
\begin{equation}
\tilde{a}_{t}^{(j)} = a_{t}^{(j)} \wedge \mathcal{R}\!\left(j,t,\mathbf{s}_t\right),
\label{eq:intro_neurosymbolic}
\end{equation}
where $a_{t}^{(j)}$ is the statistical anomaly flag for target $j$,
$\mathbf{s}_t$ denotes contextual state information, and $\mathcal{R}(\cdot)$
represents the reasoning process that validates or vetoes the candidate. This
separation between learned deviation detection and explicit rule-based refinement
provides the conceptual basis for the framework developed in this thesis.

\section{Problem Statement}

This thesis studies anomaly detection in hourly multivariate time series derived from the Austrian electricity domain. The data consist of interacting market and grid variables--including electricity price, system load, solar generation, wind generation, and derived capacity-related targets--augmented with meteorological variables, calendar descriptors, and engineered temporal features. Let $\mathbf{x}_t \in \mathbb{R}^{d}$ denote the feature vector at time $t$, and let the multivariate target vector be
\begin{equation}
\mathbf{y}_t = \bigl[y_t^{(1)}, y_t^{(2)}, \dots, y_t^{(m)}\bigr]^{\top},
\end{equation}
where, in the implemented pipeline, the principal targets correspond to solar capacity factor, wind capacity factor, load, and day-ahead price. The task is to learn normal system behaviour from historical data and then identify meaningful deviations on unseen future periods.

A forecasting-based anomaly detector is adopted because large-scale, trustworthy anomaly labels are typically unavailable in operational electricity settings. Under this formulation, the model first predicts the expected target vector using \eqref{eq:intro_forecast}, then computes residuals using \eqref{eq:intro_residual}. Because large-scale, trustworthy anomaly labels are typically unavailable in
operational electricity settings, the task is formulated through forecasting residuals.
Using the predictions and residuals defined in
\eqref{eq:intro_forecast} and \eqref{eq:intro_residual}, a target-specific
anomaly candidate is flagged when
\begin{equation}
a_{t}^{(j)} = \mathbb{I}\!\left(\left|r_{t}^{(j)}\right| > \tau_j(c_t)\right),
\label{eq:intro_threshold}
\end{equation}
where $c_t$ denotes the operating context at time $t$ and $\tau_j(c_t)$ is a
context-dependent threshold calibrated from validation residuals. Contextual
calibration is used because predictability differs across hours, seasons, and
special-calendar periods
\cite{hong2016probabilistic,weron2014electricity,shafer2008tutorial}.

The resulting decision problem is multivariate and operational: signals are
interdependent, residual variability changes across operating regimes, and some
statistical candidates require physical or market context before they can be
regarded as meaningful. The framework must also remain computationally efficient
and reproducible enough to support deployment-oriented evaluation.

To rank detected deviations, the framework additionally defines a composite severity score over all targets,
\begin{equation}
s_t = \sum_{j=1}^{m} \max\!\left(0, \frac{\left|r_t^{(j)}\right|}{\tau_j(c_t)+\varepsilon} - 1\right),
\label{eq:intro_score}
\end{equation}
where $\varepsilon > 0$ is a small constant for numerical stability. This score
supports event construction and downstream interpretation, while the symbolic
stage converts the raw candidates into a contextually refined anomaly set.

Accordingly, the problem addressed in this thesis can be stated as follows: \emph{given an hourly Austrian multivariate electricity dataset with weather and calendar context, and given a strict chronological train--validation--test protocol, design a lightweight and reproducible framework that learns normal behaviour through forecasting, detects anomalies through calibrated residual analysis, and refines those anomalies through symbolic reasoning to obtain interpretable, plausible, and decision-relevant abnormality signals.}

\section{Research Objectives \& Questions}

The overall objective of this thesis is to develop and evaluate a complete
dCeNN--ELM--ASP framework for anomaly detection in weather-dependent electricity
markets and grid operations. The work pursues the following specific objectives:
\begin{enumerate}
    \item to design a reproducible anomaly detection pipeline covering data engineering, chronological splitting, model training, threshold calibration, anomaly detection, symbolic refinement, event construction, evaluation, and deployment-oriented export;
    \item to develop a lightweight multivariate forecasting backbone based on a Tiny dCeNN encoder and an ELM-style closed-form ridge regression head, thereby targeting a favourable trade-off between predictive capability and computational efficiency \cite{chua1988cellular,huang2006extreme,hoerl1970ridge};
    \item to model uncertainty explicitly by calibrating adaptive residual thresholds on validation data, with optional contextual bucketing so that anomaly decisions reflect heterogeneous operating regimes rather than a single static tolerance \cite{shafer2008tutorial};
    \item to incorporate Answer Set Programming as a symbolic refinement layer that suppresses implausible or contextually explainable anomaly candidates and thereby improves consistency, plausibility, and interpretability \cite{gelfond1988stable,gebser2012answer,gebser2019multi};
    \item to evaluate the framework not only in terms of predictive quality, but also in terms of event-level interpretability, utility-oriented relevance, and CPU-first deployability through comparison against conventional Ridge and LSTM baselines \cite{hoerl1970ridge,hochreiter1997long}.
\end{enumerate}

These objectives lead to the following research questions:
\begin{enumerate}
    \item \textbf{RQ1: Forecasting backbone suitability.} Can a lightweight dCeNN--ELM model learn the normal multivariate dynamics of market and grid signals well enough to serve as a reliable basis for residual-based anomaly detection?

    \item \textbf{RQ2: Uncertainty-aware detection quality.} Do validation-calibrated, context-aware residual thresholds yield more stable and relevant anomaly decisions than static thresholds in a weather-driven electricity setting?

    \item \textbf{RQ3: Value of symbolic refinement.} Does the ASP-based reasoning layer reduce false positives and improve the plausibility and interpretability of detected anomalies by enforcing explicit domain constraints?

    \item \textbf{RQ4: Practical deployability.} Compared with benchmark alternatives such as Ridge regression and LSTM forecasting, does the proposed framework provide a more favourable efficiency--performance trade-off for CPU-oriented or edge-oriented deployment?
\end{enumerate}

Together, these questions define the scientific and engineering scope against which
the framework is evaluated in the later chapters.

\section{Thesis Contributions}


The contribution of this thesis lies in integrating established learning,
thresholding, and reasoning techniques into a coherent anomaly-detection framework
for weather-dependent electricity systems. The main contributions are summarized
as follows:

\begin{enumerate}
	\item \textbf{A hybrid neuro-symbolic anomaly detection architecture.}
	A unified dCeNN--ELM--ASP pipeline combines lightweight multivariate forecasting
	with rule-based validation of anomaly candidates.
	
	\item \textbf{Context-aware residual detection under weather uncertainty.}
	Validation-based, target-specific thresholds adapt anomaly decisions to
	time-varying operating conditions rather than relying on a single global
	tolerance
	\cite{hong2016probabilistic,weron2014electricity,shafer2008tutorial}.
	
	\item \textbf{Logical filtering for plausibility and interpretability.}
	ASP rules refine statistical candidates using explicit physical, temporal,
	and market constraints and provide traceable reasons for the resulting
	decisions \cite{gelfond1988stable,gebser2019multi}.
	
	\item \textbf{Event-level and utility-oriented anomaly interpretation.}
	Timestamp-level detections are organised into events and examined from an
	operational and economic perspective.
	
	\item \textbf{A reproducible and deployment-conscious implementation.}
	The work provides a modular, configuration-driven experimental pipeline with
	baseline comparison and compact model export.
\end{enumerate}

Together, these contributions define the integrated methodological and
implementation scope evaluated throughout the remainder of the thesis.

